% Point to the main file
\documentclass[../../Main.tex]{subfiles}

% ========== Start the document ==========
\begin{document}

\section{Section 1.1 - Intro to Logic}

\subsection{Propositions}

\begin{definition}[Proposition]
    A \textbf{proposition} is a statement that is either 
\end{definition}

\begin{example}[Examples with Propositions]
    Determine if the following are propositions

    \begin{enumerate}
        \item It's all coming back to me now
            \begin{solution}
                yes
            \end{solution}

        \item Hello.
            \begin{solution}
                No
            \end{solution}

        \item $x + 3 = 10$
            \begin{solution}
                No, since we don't know $x$
            \end{solution}

        \item In a right triangle, $a^2 + b^2 = c^2$
            \begin{solution}
                No, since $a, b, c$ are not assigned anything
            \end{solution}

        \item In a right traingle having legs of length $a$ and $b$, and hypotenuse of $c$, then $a^2 + b^2 = c^2$
            \begin{solution}
                Yes, since $a,b,c$ are properly specified
            \end{solution}

        \item There are 75 people in this room or the sky is purple
            \begin{solution}
                Yes
            \end{solution}
            
        \begin{note}
            7 is a \textbf{compound proposition}
        \end{note}

    \end{enumerate}
\end{example}

\begin{definition}[Compound Proposition]
    A \textbf{compound proposition} is a proposition that contains two or more \textbf{atomic} propositions. In the previous example, the two are:

    \begin{enumerate}
        \item There are 75 people in this room
        \item The sky is purple
    \end{enumerate}
\end{definition}

\begin{definition}[Atomic Proposition]
    An \textbf{atomic proposition} is a proposition that cannot be broken down further into seperate propositions
\end{definition}

\begin{note}
    We denote seperate propositions as: $p , q, r$
\end{note}

\subsection{Logical Operations (connectives)}

\begin{tablewithheader}
    { |c|c|c|c| }
    { Name & Notation & Phrasing & Truth Value}

    Disjunction (or) & $p \lor q$ & "p or q" & True when $p,q$ or both are \TRUE \\
    Conjunction (and) & $p \land q$ & "p and q" & True when both are \TRUE \\  
    Negation (not) & $\neg p$ & "not p" & The opposite truth value of $p$ \\
    Exclusive Or (XOR) & $p \oplus q$ & "p or q, but not both" & \TRUE when \textbf{only one} is \TRUE \\

\end{tablewithheader}

Here are the truth tables of each operation

\begin{tablewithheader}
    { |c|c|c|c|c|c| }
    { $p$ & $q$ & $p \lor q$ & $p \land q$ & $\neg p$ & $p \oplus q$ }

    \T & \T & \T & \T & \F & \F\\
    \T & \F & \T & \F & \F & \T \\
    \F & \T & \T & \F & \T & \T \\
    \F & \F & \F & \F & \T & \F \\

\end{tablewithheader}

\begin{example}[Logical Operators Examples]

    Let $p$ be the proposition "a is odd" and $q$ be "a is even" where $a = 3$ \\

    \begin{center}
        $p$: \; 3 is odd = \TRUE
        \\
        $q$: \; 9 is even = \FALSE
    \end{center}

    Write the following Propositions:

    \begin{enumerate}
        \item $\neg p$ \; (3 is even)
        \begin{solution}
            \FALSE
        \end{solution}

        \item $\neg q$ \; (9 is odd)
        \begin{solution}
            \TRUE
        \end{solution}

        \item $p \lor q$ \; (3 is odd or 9 is even)
        \begin{solution}
            \TRUE since 3 is odd
        \end{solution}

        \item $p \land q$ \; (3 is odd and 9 is even)
        \begin{solution}
            \FALSE since 9 is not even
        \end{solution}

        \item $p \land \neg q$ \; (3 is odd and 9 is odd)
        \begin{solution}
            \TRUE since both 3 is odd and 9 is odd
        \end{solution}    
    \end{enumerate}

\end{example}

\subsection{Conditionals}

\begin{tablewithheader}
{ |c|c|c|c| }
{ Name & Notation & Phrasing & Truth Value }
    Conditional & $p \implies q $ & "p implies q" or "if p then q" & Always True except when $T \implies F $ \\
    Biconditional & $p \iff q$ & "p if and only if q" & True when $p$ and $q$ are the same value \\
\end{tablewithheader}

Here are the truth tables

\begin{tablewithheader}
{ |c|c|c|c| }
{  $p$ & $q$ & $p \implies q$ & $p \iff q$ }
    \T & \T & \T & \T \\
    \T & \F & \F & \F \\
    \F & \T & \T & \F \\
    \F & \F & \T & \T \\
\end{tablewithheader}

\subsection{Logical Transforms}

\begin{tablewithheader}
{ |c|c|c|c| }
{ Name & Original & Output }
    Converse (swap) & $p \implies q$ & $q \implies p$ \\
    Inverse (negate) & $p \implies q$ & $\neg p \implies \neg q$ \\
    Contrapositive (converse + inverse) & $p \implies q$ & $\neg q \implies \neg p$ \\
\end{tablewithheader}

Here are the truth tables

\begin{tablewithheader}
{ |c|c|s|c|c|s| }
{ $p$ & $q$ & $p \implies q$ & $\neg q$ & $\neg p$ & $\neg q \implies \neg p$ }
    \TRUE & \TRUE & \TRUE & \FALSE & \FALSE & \TRUE \\
    \TRUE & \FALSE & \FALSE & \TRUE & \FALSE & \FALSE \\
    \FALSE & \TRUE & \TRUE & \FALSE & \TRUE & \TRUE \\
    \FALSE & \FALSE & \TRUE\ & \TRUE & \TRUE & \TRUE \\
\end{tablewithheader}

\begin{proposition}[Contrapositive Correlation]
    From the truth table, we see that the converse is logically equivalent to the original expression, so

    \begin{center}
        $p \implies q \equiv \neg q \implies \neg p$
    \end{center}
\end{proposition}

\begin{definition}[Logical Equivalence]
    Two propositions are \textbf{logically equivalent} if they have the same truth tables
\end{definition}

\begin{example}[Logical Transforms Examples]
    Let $p$ be the proposition "it's cold outside" and $q$ be "it's snowing outside"
    
    \begin{enumerate}

        \item $p \implies q$ (if it's cold outside, then it's snowing outside)
        \begin{solution}
            \FALSE
        \end{solution}
        
        \item $q \implies q$ (converse) (if it's snowing outside, then it's cold outside)
        \begin{solution}
            \TRUE
        \end{solution}
        
        \item $\neg p \implies \neg q$ (inverse) (if it's not cold outside, then it's not snowing)
        \begin{solution}
            \TRUE
        \end{solution}
        
        \item $\neg q \implies \neg p$ (contrapositive) (if it's not snowing outside, then it's not cold outside)
        \begin{solution}
            \FALSE
        \end{solution}

    \end{enumerate}
\end{example}

\begin{example}[More Examples with Logical Transforms]

    Let $p$ be the proposition "The triangle is equilateral" and $q$ be "The triangle is equiangular"

    \begin{enumerate}

        \item $p \implies q$ (if the triangle is equilateral, then triangle is equiangular")
        \begin{solution}
            \TRUE
        \end{solution}
    
        \item $q \implies p$ (converse) (if the triangle is equiangular, then the triangle is equilateral)
        \begin{solution}
            \TRUE
        \end{solution}
    
        \item $\neg p \implies \neg q$ (inverse) (if the triangle is not equiangular, then the triangle is not equilateral)
        \begin{solution}
            \TRUE
        \end{solution}
    
        \item $\neg q \implies \neg p$ (contrapositive) (if the triangle is not equiangular then the triangle is not equilateral)
        \begin{solution}
            \TRUE
        \end{solution}

    \end{enumerate}
\end{example}

\begin{example}[Truth Table Examples]

    \;

    \begin{enumerate}
        \item
            $\neg p \lor q$
            
            \begin{tablewithheader}
            { |c|c|c|c| }
            { $p$ & $q$ & $\neg p$ & $\neg p \lor q$ }
            
                \T & \T & \F & \T \\
                \T & \F & \F & \F \\
                \F & \T & \T & \T \\
                \F & \F & \T & \T \\
            
            \end{tablewithheader}

        \item
            $(p \implies r) \land (q \implies r)$

            im too lazy to fill this out
        
    \end{enumerate}
\end{example}

% ========== End the document ==========
\end{document}